\documentclass[12pt,a4paper]{article}

% Mathematics
\usepackage{amsmath,amssymb}
\usepackage{bbm}
\usepackage{slashed}

% Graphics
\usepackage{graphicx}
\usepackage{epsfig}
\usepackage{epstopdf}

% Tables and page layout
\usepackage{multirow}
\usepackage{placeins}
\setlength{\parskip}{10pt}

% Text formatting and URLs
\usepackage[normalem]{ulem}
\usepackage{url}

% Colors and cross-document citations
\usepackage[dvipsnames]{xcolor}
\usepackage{xcite}
\externalcitedocument{main-v6}

% Hyperlinks (load near the end)
\usepackage[
    colorlinks=true,
    citecolor=purple,
    linkcolor=blue
]{hyperref}

% Custom commands
\newcommand{\GJH}[1]{\textcolor{blue}{GJH: #1}}
\newcommand{\WXP}[1]{\textcolor{magenta}{#1}}

\begin{document}

%\title{Referee Report on DY13253}
%\author{Rundong Fang, Jinhui Guo, Ming Li, Jia Liu, Xiao-Ping Wang, Yiheng Xiong}
%\date{2026.19.22}
%\maketitle

\noindent \textbf{Replies to Referee}
\medskip

\subsection*{Replies to Referee}

\noindent Dear Editor,

This paper considers an anomaly-free, renormalizable model of an additional Abelian
gauge boson \(A'\) with chiral couplings to muons and taus. In this model, vector-like
leptons and a scalar field \(S\) are introduced to construct the ultraviolet completion
of the axial \(U_A(1)_{L_{\mu-\tau}}\) gauge theory. The phenomenology of the \(A'\) in
the axial limit is studied, and the relevant experimental constraints are discussed.
Additionally, possible collider signatures of the model are investigated.

The referee summarizes our work as follows:

\begin{quote}
\itshape
The authors construct an anomaly-free, renormalizable model in which an additional Abelian gauge field \(A'\) has chiral couplings to muons and taus. This model realizes a limit where the \(A'\) coupling to the \(\mu-\tau\) current is axial. The second goal of the paper is to examine the phenomenology of this model in the axial limit.

The paper presents results that are useful for the model building community. However, I have some significant reservations about key points that either should be included or otherwise justified in being excluded from the manuscript. These are explained below. Revisions for these major points (even if only to clarify why they are not issues) are required before the paper is at the level of a PRD manuscript.
\end{quote}

We sincerely thank the referee for the positive assessment of the model-building construction and for identifying issues that require further clarification or quantitative analysis. We respond to each point below.

\medskip
\noindent {\bf Major Points 1)}. \begin{quote}\itshape
Phenomenology of the \(S\) scalar. The second half of the paper reviews a set of experimental bounds on only an effective theory of the \(A'\) axially coupled to the \(L_{\mu-\tau}\) current. This leaves a major gap between these two parts of the paper: the authors show that their UV theory requires a scalar \(S\) that cannot be arbitrarily decoupled from the low-energy dynamics, e.g. from Eq. (34) which states that \(m_S^2 < 4\pi\left(m_{A'}/(2g_X)\right)^2\).\end{quote}

\noindent {\bf Reply.}
We thank the referee for highlighting the need to clarify the role of the scalar \(S\) in the UV completion and its phenomenological implications. We agree that \(S\) is an intrinsic component of the symmetry-breaking sector, and that its phenomenological relevance must be assessed under clearly stated assumptions and in light of the applicable theoretical and experimental constraints.

In the negligible-Higgs-portal-mixing benchmark adopted in the current manuscript, the scalar mass relations are
\[
m_{A'}=2g_Xv_s,\qquad m_s^2=2\lambda_Sv_s^2,
\qquad
\lambda_S=\frac{2g_X^2m_s^2}{m_{A'}^2}.
\]
Imposing the perturbativity criterion \(\lambda_S\leq\lambda_S^{\max}\) therefore gives
\[
m_s\leq\sqrt{2\lambda_S^{\max}}\,v_s
=\sqrt{\frac{\lambda_S^{\max}}{2}}\frac{m_{A'}}{g_X}.
\]

We emphasize that this theoretical upper bound does not, by itself, establish an experimentally allowed scalar-mass interval. A consistent phenomenological classification must distinguish parameter points that satisfy the adopted theoretical criterion, points excluded by the applicable experimental constraints, points not excluded by the constraints considered, and points that have not yet been assessed. The resulting scalar-mass classification is summarized as follows: \textbf{[To be completed after the scalar-sector analysis and applicable experimental recasts, specifying the excluded, not-excluded-under-the-considered-constraints, and unassessed regions.]}

\medskip
\noindent {\bf Major Points 1.1)}. \begin{quote}\itshape
It is thus an odd choice to exclude the scalar from the phenomenology of the effective theory: if the \(S\) is unavoidable at phenomenologically relevant scales, shouldn't one of the main points of the paper be that the only known UV completion of an axial \(L_{\mu-\tau}\) gauge boson has a prediction for a scalar particle in a mass range that can be probed by colliders?\end{quote}

\noindent {\bf Reply.}
We agree that the scalar prediction should be assessed within the UV completion rather than dismissed solely by invoking a scalar-decoupling benchmark. However, a scalar mass below the perturbativity upper bound does not, by itself, establish collider sensitivity. A quantitative assessment requires the production rate, physical Higgs--singlet mixing, total width, branching fractions, lifetime, and the assumptions underlying the relevant experimental analyses to be consistently matched to the model.

In particular, when kinematically allowed, the decay channel \(s\to A'A'\) can modify the visible branching fractions and thereby affect the interpretation of collider constraints. We therefore emphasize that collider sensitivity must be evaluated using the physical branching fractions and the assumptions of the applicable experimental searches; the mass relation alone is insufficient to establish either an exclusion or a discovery reach.

\textbf{[To be completed: assess the applicable scalar searches and perform the model-specific signal/decay analysis and experimental recasting, including the \(s\to A'A'\) channel whenever it is kinematically open and relevant.]}

The revised manuscript will distinguish mass ranges that have been quantitatively tested against experimental constraints from those for which only mass coverage, rather than a model-specific exclusion, can be established. The resulting collider constraints and their implications for the scalar mass range are \textbf{[to be inserted after completing the signal/decay analysis and experimental recasting]}.

\medskip
\noindent {\bf Major Points 1.2)}. \begin{quote}\itshape
What is the range of phenomenologically allowed \(S\) mass including perturbativity limits on \(\lambda\) and a motivated range of \(m_{A'}\) and \(g_X\). For example: perhaps the authors imagine a range of \(m_{A'}\) below some benchmark mass where certain experimental constraints on \(g_X\) become very strong. Then with that set of constraints (which should be stated clearly, see minor point \#1 below), one may give a range of \(m_S\) and comment on the experimental constraints.\end{quote}

\noindent {\bf Reply.}
We agree that the parameter domain and the criteria used to classify it should be specified together. To establish a well-defined phenomenologically viable mass range, we will explicitly specify the scanned ranges of \(m_{A'}\) and \(g_X\), the scalar-sector benchmark, the perturbativity criterion \(\lambda_S^{\max}\), and the set of scalar-sector constraints included in the analysis. At each parameter point, \(v_s=m_{A'}/(2g_X)\) will be derived rather than varied independently, and the scalar-mass consistency condition will be evaluated using the same set of parameter choices. \textbf{**[To be completed: specify the scan ranges, scalar benchmark, perturbativity criterion, and scalar-sector constraints.]**}

The resulting mass projection need not form a single continuous interval, since different constraints may exclude disconnected regions of parameter space. We will therefore distinguish between regions excluded by the considered constraints, regions not excluded by those constraints, and regions that have not yet been assessed. Satisfying the perturbativity criterion alone will not be taken as evidence of experimental viability. Any quoted mass range will be explicitly understood as conditional on the adopted benchmark, parameter domain, and constraints included in the analysis.

The final results will be presented in Fig. [X] / Table [X], with the benchmark assumptions, applicable constraints, and classification of the parameter regions stated explicitly. \textbf{[Numerical results to be inserted after completing the parameter scan and experimental recasting, including the resulting mass range and the boundaries of the excluded and not-excluded regions.]}

\medskip
\noindent {\bf Major Points 1.3)}. \begin{quote}\itshape
Above (34): the authors ``assume negligible Higgs-portal mixing.'' Is this interaction not generated at loop level? It seems that \(S\)-\(h\) mixing cannot be avoided due to \(\mu\) and \(\tau\) loops. In turn, this seems like it would lead to strong experimental bounds on the UV completion. In fact, once there is mixing, is there not a resonance where the mixing becomes large as \(m_S\to m_h\)? It would be useful to quantify the size of this mixing over the allowed range of \(m_S\). Existing Higgs measurements may set the strongest bounds on the UV completion. If these bounds are obviously negligible (though again, I expect this to be strongly constrained as \(m_S\to m_h\)), then the authors should present this as the justification for ignoring \(m_S\) in the low-energy theory. (Though I am still skeptical that the \(S\) can be ignored since it cannot formally decouple.)\end{quote}

\noindent {\bf Reply.}
We thank the referee for highlighting the importance of scalar mixing and Higgs-portal effects in the near-degenerate mass region. In the revised manuscript, we have explicitly defined the Higgs-portal coupling and specified its renormalization condition. We have calculated the relevant contributions from the charged-lepton and vector-like-lepton sectors, carefully distinguishing the renormalized portal parameter from the momentum-dependent off-diagonal scalar self-energy.

In the vicinity of \(m_s \simeq m_h\), we have determined the physical scalar masses and mixing using the full scalar mass matrix, incorporating loop corrections where necessary. We have not extrapolated the small-mixing approximation, which is valid away from mass degeneracy, into the near-degenerate region. We have subsequently applied the relevant Higgs signal-strength and exotic-decay constraints to the physical scalar mass eigenstates.

The resulting scalar mixing and the corresponding phenomenological constraints are quantified as follows: [insert the calculated mixing and the resulting allowed or excluded parameter regions]. The relevant calculation, numerical results, and discussion have been incorporated into the revised manuscript (Sec. [X], Eqs. [X–X], and Fig. [X], as applicable).

\medskip
\noindent {\bf Major Points 2)}. \begin{quote}\itshape
Motivating the range of effective theory parameters. To further aid in the comparison between the UV theory and the two-parameter low-energy theory, it would be useful to explicitly state in one place (perhaps in a table or a figure like Fig. 2) the range of the effective model parameters that are feasible. The main phenomenological results are in Fig. 2, but it was not clear why this range of couplings and \(A'\) masses were chosen: are values outside this range not allowed? What, if anything, does the perturbativity of the UV theory impose on the effective parameters?\end{quote}

\noindent {\bf Reply.}
We agree that the parameter domain displayed in the figure must be distinguished from a complete theoretical or experimental boundary. In the revised manuscript, we have clarified this distinction. In the originally submitted manuscript, Fig.~2 displays approximately
\[
10^{-3}\,\mathrm{GeV}\leq m_{A'}\leq10^2\,\mathrm{GeV},
\qquad
10^{-6}\leq g_X\leq10^{-1}.
\]
These ranges define the plotting window only. Parameter points outside this window are not thereby excluded, and points inside it are not automatically consistent with all theoretical requirements and experimental constraints.

We have specified the relevant parameter definitions and their interrelations. In the minimal one-VEV benchmark, the symmetry-breaking scale is determined by
\[
v_s=\frac{m_{A'}}{2g_X},
\]
while the scalar self-coupling is related to the scalar mass through
\[
\lambda_S=\frac{2g_X^2m_s^2}{m_{A'}^2}.
\]
Consequently, the perturbativity requirement correlates \(m_s\), \(m_{A'}\), and \(g_X\), and does not, by itself, establish a universal phenomenologically allowed range for the two effective parameters.

We have organized the parameter-space classification to distinguish the plotting window, the adopted benchmark and its associated parameter-space slice, theoretical consistency, experimental exclusions, regions not excluded by the constraints considered, and regions that have not been assessed. The corresponding parameter table and figure caption specify the benchmark assumptions and the scope of the constraints used. \textbf{[To be completed: insert the finalized parameter table and figure caption, including the explicit scan ranges, benchmark assumptions, perturbativity criterion, applicable experimental constraints, and the resulting classification of the excluded, not-excluded-under-the-considered-constraints, and unassessed regions.]}

\medskip
\noindent {\bf Major Points 3)}.

\begin{quote}\itshape
Are neutrino oscillations possible? The Standard Model is defined with only left-handed neutrinos because of the possibility that neutrinos are their own antiparticle. In such a case neutrino masses are generated through, e.g. through the Weinberg operator.
\end{quote}

\medskip
\noindent {\bf Reply to Point 3.1}

\begin{quote}\itshape
In this model, \(U_A(1)_{L_{\mu-\tau}}\) invariance allows Weinberg operators for each generation. However, it appears that both the Yukawa matrices and neutrino masses must be aligned and therefore a non-trivial PMNS matrix does not appear to be possible.
\end{quote}

\noindent {\bf Reply.}

We agree that the existence of gauge-invariant Weinberg operators for individual generations does not, by itself, establish compatibility with the observed non-trivial PMNS matrix. The relevant issue is whether the operator structure permitted by \(U_A(1)_{L_{\mu-\tau}}\), together with the charged-lepton Yukawa sector, allows the required mismatch between the neutrino and charged-lepton mass eigenbases. The originally submitted manuscript does not provide the operator-level analysis or explicit mass matrices needed to establish this compatibility.

\textbf{[To be completed: derive the allowed and forbidden Weinberg operators and charged-lepton Yukawa terms, construct and diagonalize the resulting mass matrices, and establish whether the adopted model permits non-trivial PMNS mixing compatible with neutrino oscillation data.]}

\medskip
\noindent {\bf Reply to Point 3.2}

\begin{quote}\itshape
Perhaps a full model would require right-handed neutrinos. These would now be charged under \(U_A(1)_{L_{\mu-\tau}}\), which likely means the second and third generation must appear as vectorlike combinations. However, it seems to me that mixing between generations is prohibited by \(U_A(1)_{L_{\mu-\tau}}\).
\end{quote}

\noindent {\bf Reply.}

We agree that introducing right-handed neutrinos alone does not establish a consistent neutrino-mass mechanism or guarantee non-trivial intergenerational mixing. Their \(U_A(1)_{L_{\mu-\tau}}\) charge assignments and the resulting allowed Dirac and Majorana mass terms determine whether the required mixing can arise. Any proposed extension must therefore be assessed at the level of its gauge-invariant interactions and resulting neutrino mass matrix, rather than assuming that the addition of right-handed neutrinos resolves the mixing problem.

\textbf{[To be completed: specify the adopted right-handed-neutrino field content and charge assignments, identify the allowed mass terms, derive and diagonalize the resulting neutrino mass matrix, and verify PMNS compatibility and consistency with anomaly cancellation and the scalar and gauge sectors, if this extension is adopted.]}

\medskip
\noindent {\bf Overall response to Major Point 3.}

\begin{quote}\itshape
Thus it is not obvious to me how neutrino oscillations can be included in this framework. If there is a simple solution, the authors should include it. If there is a known less-simple solution in the literature, the authors should cite this. Otherwise, at best this will be a confusing point. At worst, it would make the construction inconsistent with neutrino oscillation experiments.
\end{quote}

\noindent {\bf Reply.}

\textbf{[To be completed: state the adopted resolution of the neutrino-mixing issue, summarize the verified result, include relevant literature citations if applicable, and specify the corresponding modifications and section/equation references in the revised manuscript.]}

\medskip
\noindent {\bf 4)}.
\begin{quote}\itshape
Is this the full set of experimental constraints? The second half of the manuscript presents a set of experimental bounds on the two-parameter effective \(A'\) model. It is not obvious if this is the full set of constraints on the model. (See also point \#1 above: I would have liked to see a section discussing the bounds on the \(S\) within this model since it cannot be decoupled.) The following experimental bounds seem like they could be relevant. It would be helpful to either include them or justify why they are subdominant to the experiments presented.
\end{quote}

\noindent {\bf Reply.}

We agree that the constraints presented for the two-parameter effective \(A'\) model do not, by themselves, establish the completeness of the phenomenological assessment of the full UV completion. The applicability and relative importance of experimental constraints depend on the physical states involved, their interaction structures, and the assumptions defining the adopted benchmark. The constraint assessment distinguishes the scalar \(s\), the gauge boson \(A'\), finite vector-like-lepton effects, and the strict pure-axial limit. For each relevant constraint, the assessment distinguishes between constraints that have been applied, those quantitatively demonstrated to be subdominant, and those that are inapplicable for a specified physical reason.

\textbf{[To be completed: provide the finalized constraint-applicability summary, including the status and physical justification of any constraints that remain unassessed.]}

The scalar-sector constraints are addressed in connection with Major Point~1, while the four specific issues raised by the referee are discussed individually below.

\medskip
\noindent {\bf 4.1)}.
\begin{quote}\itshape
Neutrino oscillation observations. See point \#3 above. It appears that this model predicts a unit PMNS matrix, which is inconsistent with neutrino oscillations.
\end{quote}

\noindent {\bf Reply.}

We agree that compatibility with the observed neutrino oscillation data is a structural consistency requirement of the model. The minimal construction presented in the originally submitted manuscript does not provide a complete neutrino-mass mechanism or an explicit neutrino mass matrix demonstrating compatibility with the observed mass-squared differences and non-trivial PMNS mixing. This issue is addressed consistently with Major Point~3, where the neutrino-sector construction and its implications for the consistency and scope of the model are considered.

\textbf{[To be completed: state the adopted resolution and its verified results, including the neutrino mass matrix and PMNS compatibility if the neutrino sector is extended, or explicitly state the limitations of the minimal construction. Insert the corresponding section and equation references.]}

\medskip
\noindent {\bf 4.2)}.
\begin{quote}\itshape
There is no quantitative estimate of the amount of kinetic mixing generated with the ``benchmark'' choices of \(\theta\leq10^{-2}\). The authors make the insightful point that kinetic mixing goes to zero in the axial limit, but because this limit is tied to the size of the mixing (the mass of the vectorlike leptons), for the above benchmark there is non-zero kinetic mixing. It would be useful to have parametric estimates of the amount of kinetic mixing one expects so that readers may compare to experimental bounds. One could invoke parametric expressions in the literature for the relevant loop diagrams.
\end{quote}

\noindent {\bf Reply.}

We agree that imposing a vanishing kinetic-mixing coefficient at a matching scale does not, by itself, provide a quantitative estimate of the residual mixing associated with finite vector-like-lepton masses and mixing angles. In the degenerate-mass limit, the one-loop threshold contributions from vector-like leptons with opposite \(U_A(1)_{L_{\mu-\tau}}\) charges cancel under the corresponding mass and coupling assumptions. This cancellation does not establish that kinetic mixing vanishes for arbitrary finite masses, at all energy scales, or after all relevant mixing-induced contributions are included.

The distinction between the boundary condition imposed at the matching scale and the kinetic mixing relevant at experimental scales is therefore essential.

\textbf{[To be completed: specify the matching condition and renormalization prescription, provide the relevant parametric loop expression, evaluate the residual kinetic mixing for the adopted benchmark, and compare the resulting induced interactions with the applicable experimental constraints. Include the numerical estimate, its uncertainty, and the resulting phenomenological implications.]}

\medskip
\noindent {\bf 4.3)}.
\begin{quote}\itshape
Another way of reframing the above point: the bounds presented are in the pure axial limit. Since the authors present a UV model that only realizes the axial coupling in a decoupling limit, it would be useful to comment on possible experimental bounds on the small vector-like coupling in a model with non-infinite vectorlike fermion masses. If these bounds are always small compared to the axial bounds, this should be stated and justified. The importance of this point is as follows: often bounds are sensitive to the chiral structure of interactions. Sometimes a vector interaction is more tightly constrained than its axial counterpart. Since the authors focus on the phenomenology in the axial limit, one wants to make sure that any small-but-finite vector interaction isn't significant because the constraints on the coupling are orders-of-magnitude larger than the axial coupling.
\end{quote}

\noindent {\bf Reply.}

We agree that the smallness of a residual vector coupling alone does not establish its phenomenological irrelevance, since experimental sensitivities can differ substantially between vector and axial interactions. In the conventions adopted in the manuscript, the muon couplings satisfy
\[
g_V^\mu=g_X\sin^2\theta_{\mu R},\qquad
g_A^\mu=-g_X\cos^2\theta_{\mu R},
\qquad
\left|\frac{g_V^\mu}{g_A^\mu}\right|
=\tan^2\theta_{\mu R}.
\]
Thus, for \(|\theta_{\mu R}|\leq10^{-2}\), the ratio is at most of order \(10^{-4}\). However, this parametric suppression alone is insufficient to establish that the residual vector interaction is subdominant to the axial interaction in the relevant experimental constraints.

The phenomenological implications of finite vector-like-lepton mixing must therefore be assessed by comparing the constraints applicable to the vector and axial interactions separately, while accounting for kinetic-mixing effects where relevant.

\textbf{[To be completed: provide the quantitative comparison for the adopted finite-mixing benchmark, including whether the residual vector coupling or kinetic mixing modifies the excluded and not-excluded regions shown in Fig.~2. Specify the parameter ranges, the experimental constraints considered, and the evidence supporting the resulting conclusion.]}

\medskip
\noindent {\bf 4.4)}.
\begin{quote}\itshape
Are stellar cooling bounds not relevant? See e.g. 2511.09619 for a recent paper on this. If they are not, then this is certainly worth justifying.
\end{quote}

\noindent {\bf Reply.}
We thank the referee for drawing our attention to the potentially relevant stellar-cooling constraints discussed in the cited study. We agree that their applicability to the present axial-vector model requires explicit justification. In particular, the constraints on the scalar \(s\) summarized in Table~II do not, by themselves, establish whether the gauge boson \(A'\) is subject to relevant stellar-cooling constraints.

We have assessed the applicability of the cited constraints to our axial-vector benchmark, taking into account the interaction structure, the relevant mass and coupling ranges, the production mechanisms in stellar environments, and the conditions governing trapping, escape, and decay. We have also distinguished the vector \(L_\mu-L_\tau\) current considered in the cited study from the axial-vector interactions in our model, accounting for the potentially relevant effects of longitudinal modes, residual vector couplings, and kinetic mixing.

\textbf{[To be completed: insert the verified outcome of the stellar-cooling analysis, including whether the cited constraints apply to the adopted benchmark, the resulting quantitative exclusion or comparison, or the physical justification for their subdominance or inapplicability. Specify the relevant mass and coupling ranges, the assumptions supporting the conclusion, and the corresponding section, figure, or table references in the revised manuscript.]}

\medskip
\noindent {\bf Minor Point 1)}. \begin{quote}\itshape
Range of parameters. The general construction of the model is clearly written, but it is difficult for the reader to envision the mass and coupling range that the authors are targeting. It would be very helpful if the authors decide on a range of \(m_{A'}\) and \(g_X\) of interest. For example: is there a range of masses where experimental or theoretical limits are less constraining? Is there some set of phenomenology that the authors seek to realize with a ``light'' \(A'\)? If so, they should define ``light.'' (See, e.g. the first word of the introduction, which is not defined with respect to any reference mass scale.)\end{quote}

\noindent {\bf Reply.}
We agree that the target parameter domain should be specified and physically motivated, and that the term ``light'' should not be used without a clearly stated reference scale. We have therefore removed the unqualified term ``light'' from the opening sentence of the Introduction.

The parameter scan shown in Fig.~2 covers
\[
10^{-3}\,\mathrm{GeV}\leq m_{A'}\leq10^2\,\mathrm{GeV},
\qquad
10^{-6}\leq g_X\leq10^{-1}.
\]
These intervals define the plotting window used in our phenomenological analysis; they do not, by themselves, represent theoretical boundaries or experimentally allowed regions. The choice of these ranges is motivated by [specify the relevant phenomenological targets, experimental searches, and theoretical considerations supporting the selected endpoints]. We distinguish this scan window from the parameter region that remains viable after the applicable theoretical and experimental constraints are imposed.

\medskip
\noindent {\bf Minor Point 2)}.
\begin{quote}\itshape
Defining the pure-axial limit. Equation (12) gives the coupling of the \(A'\) to the mass-basis fields. It is here that the ``pure axial coupling'' limit is defined? i.e. \(\theta_{\mu R}\to0\)? If that is the case, it would help to define this as the pure-axial limit, a phrase which is used below (36) but not otherwise quantified.
\end{quote}

\noindent {\bf Reply.}

We agree that the pure-axial limit should be defined explicitly when the mass-basis couplings are introduced. In the revised manuscript, we have defined this limit as \(\theta_{\mu R}\to0\), for which the muon couplings
\[
g_V^\mu=g_X\sin^2\theta_{\mu R},
\qquad
g_A^\mu=-g_X\cos^2\theta_{\mu R}
\]
reduce to \(g_V^\mu=0\) and \(g_A^\mu=-g_X\), respectively. Thus, the coupling of \(A'\) to the muon becomes purely axial in this limit. The corresponding definition for the tau sector has been stated consistently with the adopted charge and mixing conventions.

For finite \(\theta_{\mu R}\), the interaction contains a nonzero vector component and is therefore not strictly purely axial. We have clarified this distinction in the revised manuscript and use the term ``pure-axial limit'' only for the corresponding limiting case. \textbf{[Insert the section and equation references for the revised definition.]}

\medskip
\noindent {\bf Minor Point 3)}.
\begin{quote}\itshape
Bounds on the vectorlike leptons. Because the mixing angles scale inversely with the vectorlike lepton masses, it would be helpful to present lower limits on these masses. This, in turn, gives an upper limit on the \(\theta_{(\mu,\tau)R}\) mixing angles which then gives an upper limit on the \(A'\) vectorlike coupling to \(L_{\mu-\tau}\). How do these bounds compare with the \(\theta_{(\mu,\tau)(L,R)}\) bounds in equation (27)?
\end{quote}

\noindent {\bf Reply.}

We agree that the vector-like-lepton mass constraints, the bounds on the light--heavy mixing angles, and the resulting \(A'\) couplings should be presented consistently, while distinguishing theoretical relations from direct experimental limits. In the seesaw-like regime, the charged-lepton mass relations in Eqs.~(21) and (23) give
\[
m_\ell\simeq\theta_{\ell L}\theta_{\ell R}M_{E\ell},
\qquad \ell=\mu,\tau.
\]
Combining these relations with the adopted small-mixing condition in Eq.~(27),
\[
|\theta_{\ell L}|,\;|\theta_{\ell R}|\leq10^{-2},
\]
yields the following conditional mass scales at leading order:
\[
M_{E\mu}\gtrsim
\frac{m_\mu}{(10^{-2})^2}
\simeq1.06~\mathrm{TeV},
\qquad
M_{E\tau}\gtrsim
\frac{m_\tau}{(10^{-2})^2}
\simeq17.8~\mathrm{TeV}.
\]
For the degenerate-mass benchmark \(M_{E\mu}=M_{E\tau}\equiv M_E\) adopted in Eq.~(22), the tau-sector relation gives the stronger conditional requirement \(M_E\gtrsim17.8~\mathrm{TeV}\). This is a consequence of combining the leading-order mass relations with the adopted mixing-angle condition, rather than a direct experimental lower bound. We have distinguished this conditional mass scale from the collider limits and clarified that the \(M_E>10~\mathrm{TeV}\) condition in Eq.~(22), considered by itself, is weaker than the combined requirement.

The mass relations constrain the products \(\theta_{\ell L}\theta_{\ell R}\), rather than the two mixing angles independently. Consequently, a direct-search lower bound on \(M_{E\ell}\) constrains the corresponding mixing-angle product; an individual upper bound on \(\theta_{\ell R}\) requires additional information on \(\theta_{\ell L}\) or an explicit assumption about the relative sizes of the two mixing angles. The individual mixing-angle constraints in Eq.~(27) must therefore be distinguished from mass limits inferred from the charged-lepton mass relations.

The implications for the \(A'\) interactions follow from the mass-basis couplings. In the conventions adopted in the manuscript,
\[
g_V^\mu=g_X\sin^2\theta_{\mu R},
\qquad
g_A^\mu=-g_X\cos^2\theta_{\mu R},
\qquad
\left|\frac{g_V^\mu}{g_A^\mu}\right|
=\tan^2\theta_{\mu R}.
\]
Thus, the adopted condition \(|\theta_{\mu R}|\leq10^{-2}\) implies a relative vector component of at most approximately \(10^{-4}\). This is a consequence of the mixing-angle assumption, not an independent experimental bound, and does not by itself establish the phenomenological irrelevance of the residual vector interaction.

\textbf{[To be completed: verify the conditional mass scales using the exact mass diagonalization; identify the applicable direct-search limits for the singlet vector-like leptons in this model, including their decay-mode, branching-fraction, and lifetime assumptions; compare the resulting mass and mixing constraints with Eq.~(27); and state the corresponding implications for the \(A'\) couplings. Insert the verified experimental limits, supporting references, and the relevant section, equation, or table references in the revised manuscript.]}

\medskip
\noindent {\bf Minor Point 4)}.
\begin{quote}\itshape
Conservative benchmark. Equation (27) is presented as a ``conservative benchmark.'' What does this mean? A benchmark point would be a choice of these mixing angles. A constraint is an upper bound. Eq. (27) is written as an upper bound, but is described as a benchmark. This is not simply semantic: readers would want to understand if \(10^{-2}\) is an experimental upper limit or simply a convenient choice.
\end{quote}

\noindent {\bf Reply.}
We agree that the term ``conservative benchmark'' was ambiguous because Eq.~(27) specifies upper bounds on the mixing angles rather than a particular benchmark point. In the revised manuscript, we have replaced this terminology with ``small-mixing assumption'' and clarified that
\[
|\theta_{\ell L}|,\;|\theta_{\ell R}|\leq10^{-2},
\qquad \ell=\mu,\tau,
\]
is an assumption adopted for the phenomenological analysis, rather than a universal experimental upper limit on the mixing angles. The purpose of this assumption is to specify the small-mixing regime considered in our analysis; it does not, by itself, establish that the chosen values are conservative with respect to all relevant experimental constraints.

\textbf{[To be completed: summarize the applicable direct-search and precision constraints on the mixing angles, specifying the relevant experimental assumptions and confidence levels, and compare these constraints with the adopted \(10^{-2}\) criterion. State whether the adopted assumption is consistent with the constraints considered and provide the corresponding references and section or equation numbers in the revised manuscript.]}

\medskip
\noindent {\bf Minor Point 5)}.
\begin{quote}\itshape
What is the purpose of including muon collider bounds? The second half of the paper focuses on existing experimental bounds on the effective theory. It is a strange change in tone to include hypothetical muon collider bounds for future reach. The authors do not justify why a future muon collider is the way to test this model in the future and the section feels like it is tacked on. Most of the section emphasizes that the muon collider search would be a ``standard'' follow up of existing work (Ref 83) and so this section does not seem particularly novel. The authors simply focus on the forward-backward asymmetry. Why not also include polarization asymmetries? Or would there be asymmetries in the \(ee\) versus \(\tau\tau\) final states? Given the first half of the paper, though, it seems that the real ``smoking gun'' for this UV completion is the physics of the \(S\) particle which should be tested in any Higgs-factory phase of a future lepton collider. I feel that this section either adds too little and should be downplayed as an example of future tests, or otherwise should have its own full fledged exploration (which may be a separate paper).
\end{quote}

\noindent {\bf Reply.}

We thank the referee for this constructive comment. We agree that the scope and presentation of the future-muon-collider section should be consistent with the level of analysis provided in the manuscript. In the revised manuscript, we have repositioned this section as an illustrative study of possible future tests of the effective \(A'\) interactions, rather than as a comprehensive collider forecast or a claim of projected exclusion or discovery reach.

We have clarified the motivation for considering a muon collider as a complementary probe of the muon-philic mediator \(A'\), while avoiding the implication that it is uniquely suited to testing the model or that its discovery potential has been established. We have also clarified the role of the forward--backward asymmetry as an illustrative chirality-sensitive observable that can help distinguish vector and axial-vector interactions within the assumptions of the analysis. The discussion is restricted to the results actually presented and does not imply that polarization asymmetries or comparisons among the \(ee\) and \(\tau\tau\) final states have been quantitatively investigated.

We have also acknowledged the complementary importance of scalar-sector searches in a future lepton collider's Higgs-factory program, given the role of \(S\) in the UV completion. In the absence of a dedicated model-specific sensitivity analysis, we have limited this discussion to future prospects rather than making quantitative claims about the production rates, decay signatures, or experimental reach for \(S\).

Consistent with the referee's recommendation, we have shortened and reframed the section as an outlook on possible future tests, retaining the existing forward--backward-asymmetry analysis as an illustrative example. A comprehensive study of additional asymmetries, alternative final states, and scalar-sector collider signatures is beyond the scope of the present work.

\textbf{[To be completed: insert the final section reference and identify the corresponding changes to the discussion of the muon-collider analysis and future scalar searches in the revised manuscript.]}

\medskip
\noindent {\bf Minor Point 6)}. \begin{quote}\itshape
Terminology. The following points could be improved for clarity.
\end{quote}

\medskip
\noindent {\bf Minor Point 6.1)}. \begin{quote}\itshape
The word ``light'' on page 2 has no reference scale.
\end{quote}

\noindent {\bf Reply.}
We agree that the term ``light'' should not be used without an explicit reference scale. We have therefore replaced Light neutral spin-one bosons in the opening sentence of the Introduction with Neutral spin-one bosons. The Abstract retains the explicit MeV-to-TeV mass range considered in our analysis. This revision avoids using ``light'' as an unqualified description of the entire mass range studied in this work.

\medskip
\noindent {\bf Minor Point 6.2)}. \begin{quote}\itshape
At the bottom of page 11, when referring to Table II, the authors write ``representative bounds.'' What does representative mean? Does it mean these are the most stringent bounds? Are these ``some'' bounds? Is the list ``complete'' in any sense?
\end{quote}

\noindent {\bf Reply.}
We agree that ``representative bounds'' could be interpreted as implying either a complete compilation or the most stringent constraint at every scalar mass. Table~II summarizes selected published constraints on the scalar--muon coupling \(y_{s\mu\mu}\), including bounds derived from neutron-star dynamics, fifth-force searches, neutron-star cooling, and SN~1987A cooling and diffuse gamma rays. It is not an exhaustive inventory of all available constraints, nor does it imply that the listed bounds are uniformly the strongest at every mass.

The quoted limits inherit the production, trapping, and loop-matching assumptions of the corresponding analyses. We have clarified the scope of the table to distinguish these published constraints from a complete or model-independent exclusion of the scalar parameter space.

\medskip
\noindent {\bf Minor Point 6.3)}. \begin{quote}\itshape
Between page 12--13: ``...do not provide a universal \(g_X\) bound, but they show that the intermediate scalar-mass region cannot simply be assumed unconstrained.'' What does intermediate mass mean? It is also not clear what the intended content of this statement is: are the authors presenting a bound? Why would the reader assume that a new particle in this range would be unconstrained?
\end{quote}

\noindent {\bf Reply.}
We agree that the phrase ``intermediate scalar-mass region'' was insufficiently precise and did not clearly distinguish experimental search coverage from a model-specific exclusion. The relevant mass ranges are those probed by the published searches for leptophilic scalars, rather than a separately defined universal mass interval.

BaBar and Belle searched for
\(e^+e^-\to\tau^+\tau^-s\), followed by
\(s\to e^+e^-\) or \(\mu^+\mu^-\), over approximately
\(0.04\,\mathrm{GeV}\lesssim m_s\lesssim6.5\)--\(7\,\mathrm{GeV}\).
Belle~II has also reported scalar-specific interpretations in its four-muon and other relevant resonance searches. These results provide additional information on the scalar mass range, but their application to our UV completion depends on the assumed scalar couplings and decay branching fractions.

For the mass-proportional scalar benchmark used in the Belle analysis, the published limit \(\xi_{\rm lim}=0.483\) at \(m_s=1\,\mathrm{GeV}\) implies, at 90\% confidence level,
\[
v_s>\frac{v_h}{\xi_{\rm lim}}
\simeq509\,\mathrm{GeV},
\qquad
y_{s\mu\mu}<2.1\times10^{-4},
\]
where \(v_h=246\,\mathrm{GeV}\). Using \(m_{A'}=2g_Xv_s\), this gives the conditional relation
\[
g_X<9.8\times10^{-4}
\left(\frac{m_{A'}}{\mathrm{GeV}}\right).
\]
For this scalar-mass benchmark, the direct application of the visible-decay limit requires the \(s\to A'A'\) channel to be kinematically closed, or its effect on the visible branching fractions to be accounted for. The resulting relation is therefore a benchmark-dependent constraint, not a universal exclusion for the full UV completion.

In our model, the decay \(s\to A'A'\) becomes kinematically allowed when \(m_s>2m_{A'}\). When this channel is open, it can modify the visible scalar branching fractions, so the published visible-decay limits cannot be transferred directly without accounting for this additional decay mode.

We have therefore replaced the qualitative phrase ``intermediate scalar-mass region'' with an explicit discussion of the relevant searches and their applicability. We distinguish published search coverage and benchmark-dependent scalar limits from exclusions established for the full UV model, and do not infer that a parameter region is excluded or allowed solely from its mass coverage.

\medskip
\noindent {\bf Minor Point 6.4)}. \begin{quote}\itshape
Below (34): ``we impose Eq. (34) point by point'', it's not clear what is meant by ``point.''
\end{quote}

\noindent {\bf Reply.}
A ``point'' denotes one grid point \((m_{A'},g_X)\) in the parameter scan shown in Fig.~2. At each point, the symmetry-breaking scale is derived from
\[
v_s=\frac{m_{A'}}{2g_X},
\]
rather than varied independently. For a specified scalar mass \(m_s\), the scalar self-coupling is then determined by
\[
\lambda_S
=\frac{m_s^2}{2v_s^2}
=\frac{2g_X^2m_s^2}{m_{A'}^2}.
\]
The condition
\[
\lambda_S\leq\lambda_S^{\max}
\]
is consequently checked at each scan point for the adopted scalar benchmark.

We have clarified that this point-by-point check tests the theoretical consistency of the chosen scalar benchmark. It is distinct from the application of experimental constraints, which must be assessed separately.

\medskip
\noindent {\bf Minor Point 6.5)}. \begin{quote}\itshape
Page 21, first paragraph of Section E. What is the ``parameter region of interest''? This wasn't clearly defined anywhere.
\end{quote}

\noindent {\bf Reply.}
We agree that the phrase ``parameter region of interest'' did not specify the mass coverage of the searches considered in Section~III.E. The relevant four-muon searches cover approximately
\[
\begin{aligned}
\text{BaBar}:&\quad 0.212\text{--}10\,\mathrm{GeV},\\
\text{Belle II}:&\quad 0.212\text{--}9\,\mathrm{GeV},\\
\text{CMS}:&\quad 5\text{--}70\,\mathrm{GeV}.
\end{aligned}
\]
These intervals describe the mass coverage of the published analyses; they do not imply that the entire intervals are excluded in our model.

The relevant processes are
\[
e^+e^-\to\mu^+\mu^-A',
\qquad A'\to\mu^+\mu^-,
\]
and
\[
Z\to\mu^+\mu^-A',
\qquad A'\to\mu^+\mu^-.
\]
The reinterpretation accounts for the production rate and \(\mathrm{Br}(A'\to\mu^+\mu^-)\). The experimental analyses also impose prompt-resonance selection requirements. Their limits are applicable to our model only where the predicted \(A'\) decay length is compatible with those requirements; the mass coverage alone does not establish this condition throughout the parameter space.

As an illustrative check, at \(m_{A'}=0.212\,\mathrm{GeV}\) and \(g_X=7\times10^{-4}\), the neutrino decay channels alone give \(c\tau_{A'}\lesssim0.072\,\mu\mathrm{m}\), corresponding to a laboratory-frame decay length of approximately \(1.8\,\mu\mathrm{m}\) for \(\gamma\simeq25\). This estimate supports the prompt-decay approximation at this benchmark point, but does not establish its validity throughout the full parameter space.

The vector-to-axial rate-matching approximation is controlled for \(m_{A'}\gg m_\mu\). Near the dimuon threshold, the rate-level comparison does not by itself establish a detector-level exclusion. We have therefore clarified the experimental mass coverage and the assumptions underlying the reinterpretation, without extrapolating the prompt-resonance constraints beyond the regions where their applicability is established.

\medskip
\noindent {\bf Minor Point 6.6)}. \begin{quote}\itshape
Page 21, end of section D: What's the significance of the \(v_s=2\,\mathrm{GeV}\) line? What does it mean that \(y_{s\mu\mu}\) ``becomes large'' (nonperturbative)? Why plot it? Is it a perturbativity bound?
\end{quote}

\noindent {\bf Reply.}
The line at \(v_s=2\,\mathrm{GeV}\) is a theoretical reference associated with the benchmark \(m_s=10\,\mathrm{GeV}\) and \(\lambda_S^{\max}=4\pi\). From
\[
m_s^2=2\lambda_Sv_s^2,
\]
the scalar self-coupling condition gives
\[
v_s\geq\frac{m_s}{\sqrt{2\lambda_S^{\max}}}
=\frac{10\,\mathrm{GeV}}{\sqrt{8\pi}}
\simeq1.995\,\mathrm{GeV}.
\]
Thus, for this benchmark, the line marks the approximate lower boundary on \(v_s\); values below this boundary violate the adopted scalar self-coupling criterion. The line is not an experimental exclusion and is not a perturbativity bound on \(y_{s\mu\mu}\) itself.

Using the coupling relation adopted in our manuscript,
\[
y_{s\mu\mu}=\frac{m_\mu}{v_s},
\]
we obtain
\[
y_{s\mu\mu}\simeq
\frac{0.10566\,\mathrm{GeV}}{2\,\mathrm{GeV}}
\simeq0.053
\qquad\text{at}\qquad v_s=2\,\mathrm{GeV}.
\]
A useful measure of the perturbative expansion is
\[
\frac{y_{s\mu\mu}^2}{4\pi}
\simeq2.2\times10^{-4}\ll1.
\]
Thus, the Yukawa coupling itself is well within the perturbative regime at this reference value of \(v_s\). The phrase ``becomes large'' should therefore be understood as describing the increase of \(y_{s\mu\mu}\) as \(v_s\) decreases, rather than indicating that the coupling becomes nonperturbative at \(v_s=2\,\mathrm{GeV}\).

The line at \(v_s\simeq2\,\mathrm{GeV}\) instead corresponds to the adopted scalar self-coupling criterion for \(m_s=10\,\mathrm{GeV}\) and \(\lambda_S^{\max}=4\pi\). We have clarified this distinction in the discussion and figure caption, which identify the line as a benchmark-dependent scalar-sector consistency boundary rather than a perturbativity bound on \(y_{s\mu\mu}\) or an experimental exclusion contour.

\medskip
\noindent {\large Summary:}

We thank the referee again for the careful and constructive report. We have clarified the assumptions underlying the effective-current phenomenology, the definition of the parameter domain, and the distinction between theoretical consistency conditions and experimental constraints. We have also clarified the scope of the scalar-sector phenomenology and the assumptions under which the results obtained in the pure-axial limit can be interpreted. The remaining conclusions are stated below in terms of the model construction and the parameter regions supported by the corresponding analyses.

Regarding neutrino oscillations, we agree that compatibility with the observed neutrino mass-squared differences and non-trivial PMNS mixing is a question of model consistency rather than merely presentation. The resolution adopted in the revised manuscript is as follows: \textbf{[To be completed: state the adopted approach. If the neutrino sector is extended, summarize the resulting gauge-invariant construction, neutrino mass matrix, compatibility with oscillation data, and relevant consistency checks. If the minimal construction is retained without such an extension, state explicitly which claims have been restricted and clarify the implications and limitations of the model.]}

Regarding the scalar \(S\), we agree that its physical implications cannot be dismissed solely by invoking the effective axial theory or a scalar-decoupling assumption. We have distinguished the scalar mass and self-coupling relations implied by the symmetry-breaking sector from the experimental constraints applicable to the physical scalar state. The resulting scalar-sector conclusions are \textbf{[To be completed: summarize the verified scalar-mass range, the applicable scalar-search constraints, the effect of the \(s\to A'A'\) channel when kinematically open, and the resulting exclusions or explicitly unassessed regions.]}

We have also distinguished the mass and coupling ranges covered by the experimental analyses from parameter regions that have been quantitatively excluded in the UV completion. The interpretation of the effective-current bounds is restricted to the interaction structures and assumptions for which their applicability has been established. The final constraint classification is \textbf{[To be completed: summarize the verified effects of finite vector-like-lepton masses, residual vector couplings and kinetic mixing, together with the status of the relevant additional experimental constraints and the corresponding parameter regions.]}

The referee concludes:

\begin{quote}
\itshape
In its current form, the paper is not ready for publication. The major points above must be addressed. It's entirely possible that I have misinterpreted some of these points, but if that is the case then they at least indicate ideas that must be clarified before the manuscript merits publication. I strongly encourage the authors to further address the minor points.

The paper can be a useful addition to the model building literature. However, the question of whether this class of model can be consistent with neutrino oscillations is critical. I also feel that dismissing the physics of the scalar \(S\) ignores the actual predictions of the UV model---otherwise the authors are simply justifying the phenomenology of the effective axial theory.
\end{quote}

We appreciate the referee's emphasis on the consistency of the neutrino sector and the physical implications of the scalar sector. These concerns have guided our clarification of the relationship between the UV construction and the effective-current phenomenology. The final conclusions of the revised manuscript are restricted to those supported by the adopted model construction, the calculations performed, and the experimental constraints whose applicability has been established. We have distinguished these conclusions from any further model-dependent results that have not been demonstrated.

\medskip

\noindent Sincerely yours,

\noindent Rundong Fang, Jinhui Guo, Ming Li, Jia Liu, Xiao-Ping Wang, Yiheng Xiong

\bibliographystyle{JHEP}
\bibliography{references}
\end{document}
